%% -- TEMPLATES PARA EXERCÍCIOS PROPOSTOS -- %%
\DeclareExerciseEnvironmentTemplate{question-template}{%
    \par\vspace{0.5\baselineskip}\noindent
    \textbf{\GetExerciseProperty{counter}}\quad
}{%
    \par\vspace{0.5\baselineskip}
}
%% -- TEMPLATES PARA PROBLEMAS PROPOSTOS -- %%
\DeclareExerciseEnvironmentTemplate{problem-template}{
    \par\vspace{\baselineskip}\noindent
    {\large\textbf{\GetExerciseProperty{counter}
    \IfExercisePropertySetT{subtitle}{ \GetExerciseProperty{subtitle}}}}
    \par\nopagebreak\vspace{0.3\baselineskip}\noindent
}{\par\vspace{0.5\baselineskip}}

%> Definição do tipo question
\DeclareExerciseType{question}{
    exercise-env      = question,
    solution-env      = answer,
    exercise-name     = Exercício,
    solution-name     = Resposta,
    exercise-template = question-template,
    solution-template = question-template
}
%> Definição do tipo problem
\DeclareExerciseType{problem}{
    exercise-env      = problem,
    solution-env      = hint,
    exercise-name     = Problema,
    solution-name     = Dica,
    exercise-template = problem-template,
    solution-template = problem-template
}

%> Setup do xsim
\xsimsetup{
    path = {xsim-arquivos},
    question/within = chapter,
    question/the-counter = \thechapter.\arabic{question},
    problem/within = chapter,
    problem/the-counter = \thechapter\alph{problem},
    print-solutions/headings-template = none
}

%> Ambiente de "parts" para as questões:
\newlist{partslist}{enumerate}{1}
\setlist[partslist]{label=\textbf{\alph*.}, leftmargin=*, parsep=0pt}
\newenvironment{parts}{%
    \begin{partslist}%
    \let\part\item 
}{%
    \end{partslist}
}